\vspace{3pt plus 2pt}\noindent\hspace*{1.2em}\kbvlabel{ಉದಾ}\par\vspace{1pt}

\begin{kbvderiv}
\vspace{4pt plus 3pt minus 1pt}
\begin{center}\begin{adjustbox}{max width=\textwidth}
{\renewcommand{\arraystretch}{1.22}\setlength{\tabcolsep}{4pt}
\begin{tabular}{l c l}
ಕೈಯನ್ನು + ಮುಟ್ಟಿ & {\color{kbvgrey}$=$} & \kbvresult{ಕೈಮುಟ್ಟಿ} \\
ಕಣ್ಣನ್ನು + ತೆರೆದು & {\color{kbvgrey}$=$} & \kbvresult{ಕಣ್ಣೆರೆದು}{\color{kbvgrey} (ಕಣ್ಣೆರೆದು)} \\
ಮೈಯನ್ನು + ಮುಚ್ಚಿ & {\color{kbvgrey}$=$} & \kbvresult{ಮೈಮುಚ್ಚಿ} \\
ಕಾಲನ್ನು + ತೆಗೆದು & {\color{kbvgrey}$=$} & \kbvresult{ಕಾಲ್ದೆಗೆದು}{\color{kbvgrey} (ಕಾಲ್ದೆಗೆದು)} \\
\end{tabular}}
\end{adjustbox}\end{center}
\vspace{4pt plus 3pt minus 1pt}
\end{kbvderiv}

\kbvsechead{(8)\hspace{0.45em}ಗಮಕ ಸಮಾಸ :}

\begin{kbvrule}
\noindent{\kbvsecfont ವ್ಯಾಖ್ಯೆ :}\hspace{0.5em}ಪೂರ್ವಪದವು ಸರ್ವನಾಮವೋ ಕೃದಂತವೋ ಆಗಿದ್ದು ಉತ್ತರಪದದಲ್ಲಿರುವ ನಾಮಪದದೊಂದಿಗೆ ಸೇರಿ, ಆಗುವ ಸಮಾಸಕ್ಕೆ ಗಮಕ ಸಮಾಸವೆಂದು ಹೆಸರು.
\end{kbvrule}

\vspace{3pt plus 2pt}\noindent\hspace*{1.2em}\kbvlabel{ಉದಾ}\par\vspace{1pt}

\begin{kbvderiv}
\vspace{4pt plus 3pt minus 1pt}
\begin{center}\begin{adjustbox}{max width=\textwidth}
{\renewcommand{\arraystretch}{1.22}\setlength{\tabcolsep}{4pt}
\begin{tabular}{l c l}
ಅದು + ಬೆಟ್ಟ & {\color{kbvgrey}$=$} & \kbvresult{ಆ ಬೆಟ್ಟ}{\color{kbvgrey} (ಪೂ.ಪ. -- ಸರ್ವನಾಮ)} \\
ಇವನು + ಹುಡುಗ & {\color{kbvgrey}$=$} & \kbvresult{ಈ ಹುಡುಗ}{\color{kbvgrey} (ಪೂ.ಪ. -- ಸರ್ವನಾಮ)} \\
ಅವಳು + ನಟಿ & {\color{kbvgrey}$=$} & \kbvresult{ಆ ನಟಿ}{\color{kbvgrey} (ಪೂ.ಪ. -- ಸರ್ವನಾಮ)} \\
ತಿಂದುದು + ಕೂಳು & {\color{kbvgrey}$=$} & \kbvresult{ತಿಂದಕೂಳು}{\color{kbvgrey} (ಪೂ.ಪ. -- ಕೃದಂತ)} \\
ಬೆಂದುದು + ಬೇಳೆ & {\color{kbvgrey}$=$} & \kbvresult{ಬೆಂದಬೇಳೆ}{\color{kbvgrey} (ಪೂ.ಪ. -- ಕೃದಂತ)} \\
ಮಾಡಿದುದು + ಅಡಿಗೆ & {\color{kbvgrey}$=$} & \kbvresult{ಮಾಡಿದಡಿಗೆ}{\color{kbvgrey} (ಪೂ.ಪ. -- ಕೃದಂತ)} \\
\end{tabular}}
\end{adjustbox}\end{center}
\vspace{4pt plus 3pt minus 1pt}
\end{kbvderiv}

\kbvexercise{ಅಭ್ಯಾಸ ಪತ್ರಿಕೆ ೧೩}

\vspace{2.5pt plus 3pt minus 0.5pt}\noindent\hspace*{1.2em}\kbvnum{(1)}\hspace{0.55em}ಸಂಧಿಗೂ ಸಮಾಸಕ್ಕೂ ಇರುವ ವ್ಯತ್ಯಾಸವೇನು ?\par

\vspace{2.5pt plus 3pt minus 0.5pt}\noindent\hspace*{1.2em}\kbvnum{(2)}\hspace{0.55em}“ಮರಗಾಲು” -- ಈ ಸಂಧಿಪದದಲ್ಲಿ ಸಂಧಿಯನ್ನು ಬಿಡಿಸಿ, ಸಂಧಿಯ ಹೆಸರು ಬರೆಯಿರಿ.\par

\vspace{2.5pt plus 3pt minus 0.5pt}\noindent\hspace*{1.2em}\kbvnum{(3)}\hspace{0.55em}“ಮರಗಾಲು” -- ಈ ಪದದ ವಿಗ್ರಹವಾಕ್ಯ ಬರೆದು, ಸಮಾಸದ ಹೆಸರು ತಿಳಿಸಿ.\par

\vspace{2.5pt plus 3pt minus 0.5pt}\noindent\hspace*{1.2em}\kbvnum{(4)}\hspace{0.55em}ತತ್ಪುರುಷ ಸಮಾಸವೆಂದರೇನು ? ಉದಾಹರಣೆ ಸಹಿತ ತಿಳಿಸಿ.\par

\vspace{2.5pt plus 3pt minus 0.5pt}\noindent\hspace*{1.2em}\kbvnum{(5)}\hspace{0.55em}ಕರ್ಮಧಾರಯ ಸಮಾಸವೆಂದರೇನು ? ಉದಾಹರಣೆ ಸಹಿತ ವಿವರಿಸಿರಿ.\par
